
\subsubsection{Weight-Space Learning Methods}

Set Neural Encoder (SNE) established the current cross-architecture baseline by applying hierarchical set encoding to neural network weights \citep{Andreis2023}. SNE chunks all weights into uniform 256-dimensional tokens regardless of operation type, aggregates through Set Transformers layer-wise and network-wise, and achieves $\rho$ = 0.54 on $ResNet\rightarrow ViT$ transfer. This architecture-invariant approach enables cross-family transfer but discards operation-specific signals.

Self-Attentive Neural Embeddings (SANE) demonstrated that operation-specific weight patterns improve within-architecture transfer (Schürholt et al., 2024). SANE applies spatial tokenization to convolutional weights, chunks them into 256-dimensional tokens, and uses window-based sequential encoding with positional embeddings. Contrastive learning via NTXent loss achieves +2.2\% improvement over training from scratch for ResNet-$18\rightarrow ResNet-50$ transfer. However, SANE requires shared token sizes across architectures, limiting it to same-family transfer.

Universal Neural Functionals (UNF) constructs permutation-equivariant bases for processing weight spaces of any architecture \citep{Zhou2024}. UNF enumerates valid partitions of weight symmetries and applies equivariant linear maps proven to span all S-equivariant functions. Experiments demonstrate 10-15\% improvement over non-equivariant baselines on small-scale learned optimization tasks with RNN and Transformer models. The framework provides theoretical foundations but lacks large-scale empirical validation for cross-architecture property prediction.

CAPE integrates these approaches: modular operation encoders preserve SANE's operation-specific structure for convolutions and apply UNF's equivariance for attention operations, while contrastive projection aligns heterogeneous embeddings to achieve SNE's cross-architecture capability.

\subsubsection{Cross-Architecture Alignment}

Contrastive learning has aligned fundamentally different modalities in other domains. CLIP demonstrates that InfoNCE loss aligns text and image embeddings despite their representational differences, enabling zero-shot transfer \citep{Radford2021}. CLIP trains on 400M image-text pairs, pulling matching pairs together and pushing non-matching pairs apart in a shared embedding space.

CAPE adapts this principle to weight-space learning. Models trained on the same task (ImageNet classification) share objective functions and learned representations despite architectural differences. InfoNCE loss with temperature $\tau$ = 0.07 enforces that same-task models cluster in embedding space, creating a metric where distance reflects both task similarity and architectural similarity.

\subsubsection{Computational Graph Topology}

Graph Convolutional Networks process structured data through neighborhood aggregation \citep{Kipf2017}. Applied to neural architecture computational graphs, GCNs aggregate information from connected operations to learn topology patterns. ResNet's sequential connections with skip pathways create gradient highways, while ViT's dense intra-layer attention creates all-to-all communication patterns within blocks.

Prior weight-space learning methods primarily focused on operation types with limited attention to computational graph structure. CAPE adds a 3-layer GCN that processes architecture directed acyclic graphs (nodes = operations, edges = data flow) to capture these topological patterns as a learnable residual.
